%%%%%%%%%%%%%%%%%%%%%%%%%%%%%%%%%%%%%%%%%%%%%%
%% Created by Ariel Deardorff (ariel.deardorff@ucsf.edu), Sam Teplitzky (steplitz@berkeley.edu), John Borghi (jborghi@stanford.edu)			%%
%% Version 3.0							%%	
%% September 19, 2023						%%
%%%%%%%%%%%%%%%%%%%%%%%%%%%%%%%%%%%%%%%%%%%%%%

\documentclass{article}
\usepackage{fullpage}

\renewcommand{\familydefault}{\sfdefault}
\usepackage[scaled=1]{helvet}
\usepackage[helvet]{sfmath}
\everymath={\sf}
\usepackage[textwidth=7in,textheight=10in]{geometry}
\usepackage{parskip}
\usepackage[colorinlistoftodos]{todonotes}
\usepackage[colorlinks=true, allcolors=blue]{hyperref}
\usepackage{tabularx}
\usepackage{array}
\usepackage{titling}
\setlength{\droptitle}{-1.5cm}
\usepackage{adjustbox}
\usepackage{xcolor}
\usepackage{soul}
\usepackage{dashbox}
\usepackage{booktabs}
\usepackage{comment}

\definecolor{green}{HTML}{66FF66}
\definecolor{myGreen}{HTML}{009900}

\renewcommand{\familydefault}{\sfdefault}
\renewcommand{\arraystretch}{1.5}

\title{Open Science Team Agreement}
\author{PI/Lab/Team Name}
\date{insert date of agreement}
%\company{University of California, Berkeley}
\setcounter{tocdepth}{2}


\begin{document}

%How to use this template: This template is designed to be an open science conversation starter. To use it for your team, make a copy of the project by going to Menu > Copy Project. Learn more about the topics below, modify the highlighted sections, and delete the sections that are not relevant to your research (including this one!).

\maketitle
%\tableofcontents

%Delete this section
\section*{\textcolor{gray}{\hl{How to use this template}}}
\colorbox{yellow}{\parbox{\textwidth}{
\textcolor{gray}{The Open Science Team Agreement gives researchers and other stakeholders the tools they need to understand and advocate for open science practices at a broader scale - within their laboratory, department, or the broader community.
This template is designed to be an open science conversation starter. To use it for your team, explore the topics below, modify the highlighted sections, and delete the sections that aren’t relevant to your research (including this one!).
}}}
%%%%%%%%%%%%%%%%%%%%%%%%%

\section*{\textcolor{myGreen}{Introduction}} 

\textit{Open Science} is an important aspect of conducting scientific research. This term means different things in different teams; in our team we follow these best practices: 

\subsection*{Ethical considerations} While we aspire to practice an open model of science, we respect the complex situations that may limit the full openness of our endeavors. We practice situated openness and align our open science goals with the goals of our research and research participants. This means we restrict the sharing of sensitive data, maintain the privacy of research subjects, and aim for transparency over openness. We also recognize that power and privilege can impact one’s ability to participate in open science, but we support each other in trying to build a more equitable scientific system.

\section*{\textcolor{myGreen}{Authorship + Collaboration}}
Co-authoring and collaboration are cornerstones of our scientific work. We have an inclusive model of authorship and strive to value all contributions. We have several systems in place to facilitate our work, acknowledge contributions, and expand our network to introduce diverse perspectives. 

\textbf{Persistent Identifier} (Long-lasting reference to a digital resource)
\begin{itemize}
    \item We use \href{https://orcid.org/}{ORCID} to distinguish ourselves from other researchers and manage our identities in different submission systems.
\end{itemize}
\textbf{Authorship and Author order}
\begin{itemize}
    \item We commit to conferring authorship to all who meet the criteria and to acknowledging other contributors appropriately with the \href{https://credit.niso.org/contributor-roles/}{Contributor Roles Taxonomy}. \href{https://www.icmje.org/recommendations/browse/roles-and-responsibilities/defining-the-role-of-authors-and-contributors.html}{Learn more about criteria for authorship.}
    \item We discuss author order at the outset of a project and check in throughout the writing process. \href{https://docs.google.com/document/d/1-XypwtnwT620-mnHljhOE3ETiPL92GTH0QXtmTgrIBE/edit} {Example author order checklist}.
\end{itemize}
\textbf{Inclusive Science} (Using purposefully welcoming language and practices to promote a feeling of belonging)
\begin{itemize}
    \item We practice inclusive science by thoughtfully considering our citation networks and biases and \href{https://apastyle.apa.org/style-grammar-guidelines/bias-free-language}{using bias free language}. \href{https://complexsystemsupenn.com/diversity-1}{Learn more}. 
    \item We follow a Code of Conduct that establishes positive and prohibited team behaviors. 
\end{itemize}

\section*{\textcolor{myGreen}{Articles + Research Materials}}
We make our articles and research materials as open and accessible as possible to increase the reach and impact of our research. 

\textbf{Preregistration} (Specifying your research plan in advance and registering it in a public repository. Reduces bias in hypothesis-testing research.
\href{http://www.psychologicalscience.org/observer/research-preregistration-101#.WR3GyFPyvOT}{Learn more.})

\begin{itemize}
\item We preregister our hypothesis-testing studies in \hl{\hbox{
[\href{https://osf.io/}{Open Science Framework}/\href{https://aspredicted.org/}{AsPredicted}/other]
}}

\end{itemize}
\textbf{Methods and Protocols} (Step-by-step documents describing exactly how research was performed. Sharing methods and protocols enables other researchers to reproduce experiments.)   

\begin{itemize}
\item We publish our methods and protocols in \hl{\hbox{[\href{https://www.protocols.io/}{Protocols.io}/other]}} when the corresponding paper has been \hl{[submitted/accepted]} 
\end{itemize}

\textbf{Preprints} (Version of a paper made public prior to peer review. Sharing protocols increases the speed of research dissemination. \href{https://plos.org/open-science/preprints/}{Learn more}.) 

\begin{itemize}
\item We submit preprints of our articles to \hl{\hbox{ [\href{https://www.biorxiv.org/}{bioRxiv}/\href{https://www.medrxiv.org/}{medRxiv}/other]}} 
\end{itemize}

\textbf{Open Access} (A publishing model where articles are published online with no access restrictions so that anyone can read them)

\begin{itemize}
\item We make all of our articles openly accessible either through publishing in open access journals, or by archiving a copy in an open repository like \hl{\hbox{[our institutional repository/\href{https://www.ncbi.nlm.nih.gov/pmc/}{Pubmed Central}/other]}}
\end{itemize}

\textbf{Theses and Dissertations}

\begin{itemize}
\item Whenever possible we incorporate open science practices into the thesis writing process and try to archive a copy of team theses in our institutional repository.
\end{itemize}

\textbf{Presentations}
\begin{itemize}
\item We make our presentation slides and posters available in \hl{\hbox{[our institutional repository/\href{https://zenodo.org/}{Zenodo}/other]}} so that they are more easily discoverable and citable. 
\end{itemize}


\section*{\textcolor{myGreen}{Data + Code}}
Research data include “raw” data, processed data, data at intermediate stages, and “final” datasets (i.e. the dataset that underlies a manuscript) as well as any documentation that is needed to make use of these materials. We share our research data and code in public repositories whenever possible.

\textbf{Documentation}
\begin{itemize}
    \item We create readme documents (or equivalent) to track the data we are creating, the software we are using (including versions) and describe the code we are writing ourselves.  
\end{itemize}

\textbf{Data} 
\begin{itemize}
\item We use the \href{https://datadryad.org/stash}{Dryad Data Repository}/other to make data available to others. \href{https://www.nature.com/sdata/policies/repositories}{Find data repositories for your research}
\end{itemize}

\textbf{Software and Code} (Broadly refers to computer programs, packages, and scripts used to work with, analyze, and visualize data.)  
\begin{itemize}
\item We use \hl{\hbox{[\href{https://github.com/}{Github}/other]}} for storing code we are writing ourselves and \hl{\hbox{\href{https://zenodo.org/}{Zenodo}/other]}} for ensuring it is preserved in a citable form at the conclusion of a project. 
\item Our research relies on open-source and open infrastructure projects. We support these communities with labor, donations, and citations, and understand that our contributions are welcome and appreciated in these spaces. 
\item We assign a \hl{[MIT/Mozilla/Apache/Other]} open-source license to our code so others know how they can re-use it. \href{https://choosealicense.com/licenses/}{Learn more.}
\end{itemize}

\section*{\textcolor{myGreen}{Communication + Impact}}
Spreading the word about our research is essential for ensuring it reaches a wide audience and can have an impact. We use a variety of tools and platforms to share our research.\\

\textbf{Copyright}
\begin{itemize}
    \item When publishing and sharing research materials we aim to use open licenses, such as a creative commons licenses, to allow wide re-use of our work. \href{https://creativecommons.org/choose/}{Learn more about creative commons.}
\end{itemize}
\textbf{Research Profiles} (useful to establish a public scientific persona associated with one’s institution, co-authors and larger discipline)
\begin{itemize}
    \item We create public profiles using \hl{\hbox{[Google Scholar/University System/Other]}} to track our published or shared work.
\end{itemize}

\textbf{Social Media}
\begin{itemize}
\item We use \hl{\hbox{[University Communications Office/Twitter[X]/Discord/other platform]}} to communicate our research findings and conference presentations to a broader audience.
\end{itemize}

\section*{Keeping Ourselves Accountable}
\begin{itemize}
    \item We review this document \hl{[when a new member joins the lab, annually, at every lab meeting]}.
\end{itemize}

{\let\thefootnote\relax\footnote{{This work is licensed under a Creative Commons Attribution 4.0 International License.}}}

\end{document}